\documentclass[10pt]{article}
\usepackage[margin=20mm]{geometry}
\usepackage{amsmath,amssymb,mathtools,fontspec,unicode-math,xurl,hyperref}
\IfFontExistsTF{DejaVu Serif}{\setmainfont{DejaVu Serif}}{\setmainfont{Latin Modern Roman}}
\setmathfont{latinmodern-math.otf}
\setlength{\parindent}{0pt}\setlength{\parskip}{6pt}\sloppy
\title{Initial scientific review of
Finite index rigidity of hyperbolic groups}
\author{}
\date{Frozen batch M25-0245; initial review, decision pending}
\begin{document}
\maketitle
Abstract: The initial assessment is Minor. The low-dimensional complexity comparison, the graph-action conventions, and a zero-power boundary case require correction. The application of Proposition 7.1 also omits the standard reduction from the full Rips complex to its \(m\)-skeleton. That target choice satisfies the printed hypotheses and preserves the controlled map and the estimate needed later. The full-target compactification argument supplies additional counterevidence. The checks do not establish a counterexample to Theorems A, B, or C. The Minor findings are confined to the stated definition and convention boundaries, including this omitted standard target specification. Scientific decision and adjudication remain pending, with no human reviews in this package.\par
Article. Nir Lazarovich. Journal identifier: 10.4171/JEMS/1618.\par
Version examined. The 29-page author manuscript `arXiv:2302.04484v3`, submitted 13 October 2024. PDF and printed page numbers coincide. The frozen journal record belongs to 2025; equivalence of this manuscript with the journal version has not been established.\par
\section*{Scope and principal dependencies}
The reviewer read all 29 pages of the acquired source text and inspected the actual images of all 26 body pages. The review traced the definitions and quantifiers, the bicombing coefficient estimates, the weighted-pattern construction, the shortening argument, the compactly supported cohomology obstruction, and the passage from complexity to index rigidity. Roles and self-reported model labels are recorded separately from the controller’s UI observation.\par
Theorem B is the central estimate: for a fixed locally finite one-ended hyperbolic graph \(X\) and a free cocompact action, \[\operatorname{Vol}(X/G)\asymp_X C_{2,m}(G),\qquad m=\dim(\partial X)+1.\] The upper bound constructs a bounded-diameter subcomplex of an iterated join. For the lower bound, Proposition 6.1 controls total track weight by complexity; Proposition 7.1 and Corollary 3.7 force a positive weight near each member of a packing in the quotient. Theorem C then compares complexity to subgroup index, and the normalized limit inferior gives exact multiplicativity.\par
The construction of Mineyev’s bicombing, the full Bestvina–Mess/Walsh boundary-cohomology inputs, the cited controlled-filling results, and Sykiotis’s reduction for groups with more than one end have not been independently reproved. The full source was parsed, but that is not a certification of every external dependency. An official publisher and manuscript-record search yielded no additional source used for this review; missing code or supplementary material is not classified as a mathematical defect.\par
\section*{Confirmed issues and their limits}
\subsection*{F1. The definition at \(m=1\) does not support the stated general rank comparison (Minor)}
Location. The definition on p. 2, the comparison with rank on p. 3, and the action-based restatement and finiteness assertion on p. 25.\par
The definition on p. 2 minimizes over simplicial complexes \(\overline L\) satisfying \(\pi_1(\overline L)=G\) and having \((m-1)\)-connected universal cover. If \(m=1\), such a complex is a graph, whose fundamental group is free. Take \[G=\pi_1(\Sigma_2).\] Then \(\operatorname{rank}(G)=4\), but \(H_2(G;\mathbb Z)=\mathbb Z\), whereas a free group has zero second homology. Thus there is no admissible graph and the printed definition gives \[C_{1,1}(G)=\infty.\] This is a one-ended hyperbolic example, so the unrestricted rank comparison cannot follow merely from hyperbolicity.\par
For a nontrivial free group \(F_d\), the intended comparison does hold. A graph formed from \(d\) triangles sharing one vertex has \(1+2d\) vertices and \(3d\) edges, and hence \[d\le C_{1,1}(F_d)\le 1+5d\le 6d .\] At the trivial group, comparison with \(1+\operatorname{rank}(G)\) is the appropriate formulation. The action-based definition on p. 25 differs at \(m=1\): a connected complex with a free \(G\)-action need not be simply connected, and therefore need not be the universal cover of its quotient.\par
Impact and repair. Restrict the introductory comparison to nontrivial free groups, or define a separate action-based invariant for \(m=1\). For \(m\ge2\), the connectivity condition makes the two definitions agree. The principal estimate uses \(C_{2,m}\) with \(m\ge2\); this counterexample does not contradict that estimate.\par
\subsection*{F2. The simplicial quotient and Cayley-action conventions need a precise model (Minor)}
Location. The proof that Theorem B implies Theorem C on p. 4; the definition of a free action on p. 6 and the subdivision assertion continuing on p. 7.\par
There are two distinct points. First, one barycentric subdivision need not make the quotient an abstract simplicial complex. Let \(\mathbb Z\) act by translation on the integer line. The action is free on the realization. After one subdivision, the quotient has two vertices joined by two distinct edges. A simplicial complex cannot have two different edges on the same pair of vertices. A second subdivision gives a four-cycle. The same generator-loop phenomenon can occur in a one-ended torsion-free surface-group Cayley graph.\par
Second, vertex-freeness of an ordinary undirected Cayley action is insufficient for the geometric freeness stipulated on p. 6. Let \[G=\pi_1(\Sigma_2)\times C_2\] and include the nontrivial involution \(s\) in the generating set. The group is one-ended and hyperbolic. Left multiplication by \(s\) reverses the edge \(\{1,s\}\) and fixes its midpoint. Subdivision preserves that fixed point.\par
Counterevidence and repair. Neither observation is a counterexample to index rigidity. A uniform replacement gives the exact hypothesis needed on p. 4. Choose a finite genuine simplicial presentation complex \(P\) with \(\pi_1(P)=G\), and take \(X'\) to be the one-skeleton of its universal cover. The deck action is free and cocompact, \(X'/G=P^{(1)}\) is simplicial, and \(X'\) is quasi-isometric to \(G\), hence one-ended and hyperbolic. For every finite-index subgroup \(H\), \[\operatorname{Vol}(X'/H)=[G:H]\operatorname{Vol}(X'/G).\] Theorem B therefore gives precisely Theorem C, with constants depending on the fixed model. This also avoids claiming that subdivision can remove torsion fixed points. The later Rips target is expressly allowed to have a proper action; the finite-stabilizer extension of the pattern argument is discussed below.\par
\subsection*{F3. The proof of Corollary 1.1 omits the zero-power case (Minor)}
Location. Page 2, the assertion that the subgroup has index \(|k|\) for every \(k\).\par
Under the literal cyclic-subgroup notation, \(\langle\phi^0\rangle\) is trivial. Consequently \(F_d\rtimes\langle\phi^0\rangle=F_d\) is the kernel of the homomorphism from the mapping torus to \(\mathbb Z\), and has infinite index, not index zero. The finite-index proof applies to \(k\ne0\).\par
This does not supply a counterexample to the corollary in its usual nontrivial atoroidal setting. For \(d\ge2\), a nonzero-power mapping torus has Euler characteristic zero, while \(\chi(F_d)=1-d\ne0\). Thus one zero exponent and one nonzero exponent cannot give isomorphic groups; two zero exponents cause no difficulty. For nonzero opposite exponents, inverting the stable letter gives the expected isomorphism.\par
Repair boundary. Either specify nonzero integer exponents, or add this zero-case argument. If instead an abstract infinite cyclic factor is intended even for the identity automorphism, the notation and proof must say so; it is then the identity mapping torus \(F_d\times\mathbb Z\), not the stated index-zero subgroup.\par
\subsection*{F4. An omitted \(m\)-skeleton reduction for the Rips target (Minor)}
Location. Proposition 7.1 on p. 23 assumes that \(Y\) is \(m\)-dimensional, where \(m=\dim\partial Y+1\). Section 8 takes the full contractible Rips complex on p. 25 and applies the proposition on p. 26 to obtain (8.3).\par
The dimension need not agree. For a closed surface group, \(m=2\). In a Rips complex of sufficiently large radius \(d\ge3\), four consecutive vertices of a geodesic have pairwise distances at most three and span a \(3\)-simplex. The full Rips complex can therefore have dimension greater than \(m\). This is a counterexample to the required dimensional inference, not to quasi-surjectivity.\par
Standard target interpretation. Put \(Z=Y^{(m)}\), where \(Y\) is the full Rips complex and \(m\ge2\). The complex \(Z\) has the same one-skeleton and boundary as \(Y\), and inherits its proper cocompact action. Since \(Y\) is contractible, its \(m\)-skeleton is \((m-1)\)-connected. It has dimension \(m\). The extension from the free \(m\)-dimensional source \(L\) may be chosen cellular and equivariant, relative to the already selected geodesic edges. Its image then lies in \(Z\). Controlled relative equivariant cellular approximation in uniformly bounded simplicial stars preserves the controlled-filling inclusion (8.2), up to a constant depending only on the fixed, locally finite, finite-dimensional Rips complex. Any auxiliary subdivision is used only to construct the map and does not alter the original complexity or pattern count. Proposition 7.1 now applies with its printed hypotheses to \(Z\), and its conclusion on the common one-skeleton gives the estimate needed in (8.3). The upper-bound argument on p. 25 already uses an explicit \(m\)-skeleton. Thus the finding is the omitted skeletal reduction and literal target specification, not an established new obstruction in the core cohomology argument.\par
Additional counterevidence for the full target. The proof of Proposition 7.1 only needs the target vanishing \[\widetilde{\check H}^{\,m-1}(Y\cup\partial Y;\mathbb Z)=0.\] For the actual full Rips target, its compactification is contractible by the same input invoked on p. 24, so this vanishing holds without the target-dimension assumption. Retain the \(m\)-dimensional, \((m-1)\)-connected source \(L\). Its connecting map remains injective by Lemma 7.2. Choose \[0\ne a\in\widetilde{\check H}^{\,m-1}(\partial Y;\mathbb Z).\] Exactness and naturality give \[\Phi^*D_Y(a)=D_L((\partial\Phi)^*a)\ne0,\] because \(\partial\Phi\) is a homeomorphism and both relevant connecting maps are injective. The compactly supported class \(D_Y(a)\) is represented in \(H^m(Y,Y\setminus K;\mathbb Z)\) for a fixed compact \(K\). If \(\Phi(L)\) avoided \(K\), its pullback would be zero, a contradiction. Finally, cocompactness of \(\operatorname{Aut}(Y)\) moves this fixed compact obstruction to uniformly bounded distance from every target point. One applies the same argument after conjugating the target action and composing \(\Phi\) with an automorphism. The resulting neighborhood radius depends only on \(Y\).\par
Thus the displayed vanishing also suffices if one retains the full target. The source hypothesis and the quantitative conclusion needed for (8.3) remain unchanged. The Minor classification is based on the available standard skeletal interpretation of the actual construction; it does not reclassify a demonstrated unclosed core gap merely because a replacement proof can be supplied.\par
\section*{Substantive checks supporting the main argument}
\subsection*{The normalized limit inferior}
On pp. 4–5, order finite-index subgroups by reverse inclusion. For fixed finite-index \(H\le G\), the subgroups \(K\le H\) form the final segment starting at \(H\), so restricting to them does not change the net limit inferior. Since \[[G:K]=[G:H][H:K],\] the normalization gives \[\overline C_{2,m}(H)=[G:H]\overline C_{2,m}(G).\] The two-sided bounds in Theorem C imply \(0<\overline C_{2,m}(G)<\infty\). An isomorphism preserves both the complexity and the finite-index subgroup lattice. Equal complexities of isomorphic subgroups therefore force equal indices. Residual finiteness is not needed.\par
\subsection*{Bicombing coefficients and weighted patterns}
The checks of Lemmas 3.5–3.6 and Corollary 3.7 use the stated rational chains, without assuming a common denominator. In Lemma 3.6, the decomposition \(a=a_x+a_z+a_y\) yields boundary chains \(b_x,b_y\) with disjoint supports and augmentation one. Therefore \[\|a_z\|_1\ge\tfrac12\|\partial a_z\|_1
 =\tfrac12(\|b_x\|_1+\|b_y\|_1)\ge1.\] The fixed finite number of edges in a ball supplies the uniform coefficient bound, including endpoint choices of the point on the geodesic.\par
For the perfection step, subtracting total track defect from its maximal connector weight gives a nonnegative weight no greater than every surviving connector weight. Singular tracks then have zero surviving weight. In the resolution, a track has a single unoriented target-edge label; a domain edge has at most one connector with that label. This controls embeddedness and the crossing estimates. The triangle defect is bounded by the bicombing defect, together with a bounded central region for labels occurring on all three sides.\par
In Proposition 6.1, total edge displacement takes nonnegative integer values, so a minimizing equivariant vertex map exists. Compact regular tracks, one-endedness, and the finite number of track orbits give the bounded-side and outermost choices. Summing the shortening inequality and using minimality bounds the perfected total weight. Additive endpoint constants are absorbed by the stated \(X\)-dependent constants; they do not furnish a counterexample to the final estimate.\par
\subsection*{Proper targets and the volume lower bound}
Although Section 5 initially uses a free target, Proposition 6.1 states a proper target action. For a fixed target edge \(f\), its stabilizer is finite. The connector set has finitely many \(G\)-orbits; each such orbit contributes at most \(|\operatorname{Stab}(f)|\) connectors to the fiber of the equivariant label map over \(f\). Thus tracks remain compact, with finite stabilizer. The equivariant choices use freeness of the source \(K\), which is retained. No additional stabilizer factor is needed in the local coefficient and edge-count bounds.\par
On p. 26, the packing argument counts tracks rather than merely connectors: disjoint selected quotient balls have different target-edge orbits, and a track has only one such label. Hence their positive weights belong to distinct quotient tracks. A collapsed image edge causes no obstruction: follow a domain path to the first noncollapsed edge; a globally constant equivariant map is impossible for this proper action of an infinite group. These checks support the deduction from local bicombing mass to the lower bound in Theorem B, subject to the explicitly stated external inputs and the correction in F4.\par
\section*{Disposition}
Revise the four specified points, retaining the full group range of Theorems A–C through the precise graph conventions and the standard skeletal target specification. The initial review has not established a false main theorem. This report is a scoped scientific assessment of the acquired manuscript, not an assertion about an unexamined journal version or an independent acceptance.\par
9 N. Lazarovich, Finite index rigidity of hyperbolic groups, arXiv:2302.04484v3, 13 October 2024. Acquired source: 29 pages. European Mathematical Society Press, article record, DOI 10.4171/JEMS/1618, https://ems.press/journals/jems/articles/14298660. Used for identity and version separation. A. Dranishnikov, On Bestvina–Mess Formula, arXiv:math/0503018. Bounded primary-source cross-check of the boundary-dimension/cohomology input; not a full external proof audit.\par
\end{document}
